\section{Besoins non fonctionnels}

$ \ast $ \textit{La section suivante pourra être négociée de manière à ce que toutes les parties prenantes soient en accord avec le contenu.} \\

En plus des fonctionnalités attendues, le système devra respecter plusieurs contraintes techniques et qualitatives :


\small
\begin{itemize}
   \item \textbf{Fiabilité} : le système doit fonctionner de manière continue dans des conditions extérieures variables (pluie, vent, variations de température).
   \item \textbf{Robustesse} : les composants matériels doivent être résistants aux intempéries et aux chocs.
   \item \textbf{Évolutivité} : l’architecture logicielle et matérielle doit permettre l’ajout de nouveaux capteurs ou modules sans modification majeure.
   \item \textbf{Performance} : la détection et la transmission doivent être sans perte d’informations.
   \item \textbf{Energie et autonomie} : Comme on peut le constater dans le tableau des spécifications ci-dessus, le système doit être conçu pour une faible consommation énergétique et disposer d’une autonomie suffisante afin d’assurer un fonctionnement prolongé. Une autonomie de 24 heures a été définie afin de garantir la flexibilité du système et d’assurer sa continuité de fonctionnement.
\end{itemize}